\section{System Design}
\label{sec:design}

\system selects which Agents' histories deserve protection (O1),
which KV regions within those histories should be retained (O2), and
how much HBM replication may occupy as foreground demand changes (O3). 
Figure~\ref{fig:system-overview} shows our four cooperating components.
\emph{Tracker} estimates remaining workflow work and recovery urgency.
\emph{Router} chooses execution workers and replica destinations.
\emph{KV Master} selects Agent histories and KV regions, evaluates
candidate placements, and manages replica retention and reclamation.
\emph{Adapter} adjusts the HBM replication budget to foreground demand. Rank-local
Connectors execute transfers and report completion to KV Master.

\begin{figure}[t]
 \centering
 \includegraphics[width=\columnwidth]{figures/system-overview.pdf}
 \caption{\system: Four-component architecture. Each workflow node is one
 invocation. Letters identify Agents, and parentheses give invocation
 ordinals, B and C are each invoked twice.
 Each host contains one TP=2 DP worker. Each rank's HBM holds execution
 and replica regions. Rank-matched Connectors copy KV between hosts;
 each host's DRAM connects to its own Connectors. Blue paths carry
 KV; dashed paths carry control.}
 \label{fig:system-overview}
\end{figure}


\subsection{Tracker}
\label{sec:tracker}

% 中文：一个 Agent 通过多步完成其任务，每步对应一次模型调用。Tracker 按调用粒度
% 记录执行和历史使用关系，按 Agent 剩余工作的粒度估计其对工作流完成时间的影响。
An Agent completes its assigned task through a sequence of steps, each
corresponding to one model invocation. Tracker records execution and
history use at invocation granularity, and estimates remaining Agent
work to determine workflow urgency.

% 中文：本文考虑 Agent 任务 DAG 预先已知且固定的工作流。Agent 集合及依赖关系已知，
% 但每个 Agent 内部的调用次数和完成时间未知。
We target workflows with a fixed Agent-task DAG known before execution.
Agent identities and dependencies are fixed, while each Agent's
invocation count and completion time are unknown.

\subsubsection{Remaining-Time Prediction}
\label{sec:tracker-prediction}

% 中文：以 (w,a,r) 标识工作流 w 中 Agent a 的第 r 次调用，当依赖满足后，workflow runtime 将这次 invocation 提交给 serving 系统调度执行。
% 后续调用可以复用兼容前缀的已选择 KV；调用结束不等于历史失效，直到 Agent 结束。
Tracker identifies invocation $r$ of Agent $a$ in workflow $w$ as
$v=(w,a,r)$. Once its dependencies are satisfied, the workflow runtime
submits the invocation to the serving system for scheduling. Successive
invocations of the same Agent may reuse selected KV for a compatible
history prefix. Completing an invocation does not end the Agent's
history; Tracker tracks that history until the Agent finishes its task.

% 中文：每次调用提交时，Tracker 根据本次调用的信息预测 Agent 距离完成任务还需多久。
% 一次调用只预测一次，下一次调用到来时再更新，直到 Agent 结束。
At each invocation submission, Tracker predicts how long the Agent will
take to finish its task using the invocation's input and the Agent's
observed history. Each invocation triggers one prediction, which is
updated at each subsequent invocation until the Agent finishes its task.

% 中文：将预测建模为基于上下文的回归问题。文本编码器提取任务描述、本次新增输入
% 和近期已接受输出的语义表示，再与提交时已知的模型、角色、输入长度、已完成调用
% 次数、历史调用的长度和耗时，以及 Agent 已运行时间结合。
% 回归器直接预测 Agent 的剩余完成时间，供后续工作流 Slack 计算使用。
We formulate this prediction as context-conditioned regression. A text
encoder represents the Agent's task description, newly appended input,
and recent accepted outputs. A regressor combines this semantic
representation with features available at submission: the model and
Agent role, input length, completed-call count, previous calls' lengths
and timings, and elapsed Agent time. It directly predicts Agent $a$'s
remaining completion time, $\widehat R_{a,r}$, for workflow slack
computation.

% 中文：训练样本将每次调用提交时的特征与该次提交到 Agent 任务结束的实际时间配对。
Training examples pair these submission-time features with the observed
duration from submission of invocation $r$ to completion of Agent $a$'s
task.

% 中文：尚未启动的 Agent 使用同模型、同角色历史任务的时长先验 D。
% 首次调用提交时，以该先验初始化剩余时间估计 R，并记录提交时间 tau。
% 每次调用触发的预测返回后，以新预测替换 R，tau 始终取该预测对应的调用提交时间。
For an unstarted Agent, Tracker uses a duration prior $\bar D_a$ from
historical tasks with the same model and role. Its first invocation
submission initializes $\widehat R_a$ to this prior and $\tau_a$ to the
submission time. Each prediction $\widehat R_{a,r}$ replaces
$\widehat R_a$ when available, with $\tau_a$ set to invocation $r$'s
submission time.

\subsubsection{Slack Computation}
\label{sec:tracker-slack}

% 中文：注册工作流时，runtime 从固定编排图导出所有 Agent 任务的工作流 ID、Agent ID
% 和前驱 Agent ID 列表；边 a→b 表示 b 等待 a 完成任务后才能开始。
% runtime 报告任务完成事件，Tracker 移除已完成节点，得到剩余 DAG。
% 该 DAG 每个节点是一个尚未完成的 Agent 任务，其内部可以有多次 invocation。
At workflow registration, the runtime exports each Agent task's workflow
ID, Agent ID, and predecessor Agent IDs from the fixed orchestration graph.
An edge $a\to b$ means that $b$ waits for Agent $a$ to finish its task.
The runtime reports task completion, allowing Tracker to remove finished
tasks and maintain the residual DAG $G_t=(A_t,E_t)$. Each node in this
graph is one unfinished Agent task, which may contain multiple invocations.

% 中文：算法以剩余 Agent DAG、时长先验、已启动 Agent 的最新估计及提交时间、误差裕量为输入。
% 未启动 Agent 保留完整时长先验；已启动 Agent 的估计换算到当前时刻。
% 正向计算最早完成时间，反向计算最晚完成时间。
% 后向遍历以预测的工作流完成时间为终点；每个 Agent 的剩余工作只计算一次。
% pred 和 succ 分别表示剩余前驱与后继；g 是经校准的 Slack 高估误差裕量。
Algorithm~\ref{alg:slack-estimation} computes guarded Agent slack using
forward and backward critical-path passes~\cite{kelley1959cpm}.
Unstarted Agents retain their full duration priors, while estimates for
started Agents are adjusted to the current time $t$. The algorithm counts
each Agent's remaining work once and uses the predicted workflow finish
$T$ as the backward-pass endpoint. The two passes compute earliest and
latest finishes, $EF_a$ and $LF_a$. Here $\operatorname{pred}(a)$ and
$\operatorname{succ}(a)$ denote residual predecessors and successors,
and $g_a$ is a calibrated upper quantile of slack overestimation.

% 中文：所需时长估计缺失、过期或无效时，整个工作流暂不授予正 Slack。
% 正向遍历：EF = 本节点剩余时间 + 最晚完成的前驱；后向遍历：
% LF = 后继最晚开始时间的最小值，汇点从 T 开始回推。
% 最终 Slack 为 max(0, LF - EF - g)；两次遍历的复杂度均为 O(|A_t|+|E_t|)。
\par\medskip\noindent\begin{minipage}{\columnwidth}
\raggedright
\refstepcounter{designalgorithm}\label{alg:slack-estimation}
\textbf{Algorithm \thedesignalgorithm: Compute Agent slack}
\par\smallskip
\begin{tabularx}{\columnwidth}{@{}r@{\quad}X@{}}
\toprule
 & \textbf{Input:} $G_t=(A_t,E_t)$; Agent states; priors $\bar D_a$;
 estimates $(\widehat R_a,\tau_a)$; guards $g_a$; time $t$ \\
 & \textbf{Output:} guarded slack $\boldsymbol{s}=\{s_a:a\in A_t\}$ \\
\midrule
1 & $s_a\gets0,\quad \forall a\in A_t$ \\
2 & \textbf{for} $a\in A_t$ \textbf{do} \\
3 & \quad \textbf{if} $a$ is unstarted \textbf{then} $d_a\gets\bar D_a$ \\
4 & \quad \textbf{else} $d_a\gets\widehat R_a-(t-\tau_a)$ \\
5 & \textbf{if} any required estimate is missing or invalid, or
 any $d_a\le0$, \textbf{then return} $\boldsymbol{s}$ \\
6 & $\pi\gets\textsc{TopologicalOrder}(G_t)$ \\
7 & \textbf{for} $a$ in $\pi$ \textbf{do} \\
8 & \quad $EF_a\gets d_a+
 \max\bigl(\{0\}\cup\{EF_p:p\in\operatorname{pred}(a)\}\bigr)$ \\
9 & $T\gets\max\bigl(\{0\}\cup\{EF_a:a\in A_t\}\bigr)$ \\
10 & \textbf{for} $a$ in reverse $\pi$ \textbf{do} \\
11 & \quad $LF_a\gets
 \min\bigl(\{T\}\cup\{LF_b-d_b:b\in\operatorname{succ}(a)\}\bigr)$ \\
12 & \quad $s_a\gets\max\{0,LF_a-EF_a-g_a\}$ \\
13 & \textbf{return} $\boldsymbol{s}$ \\
\bottomrule
\end{tabularx}
\end{minipage}\par\medskip

% 中文：调用继承所属 Agent 的 Slack。
Invocation $v=(w,a,r)$ inherits its Agent's slack, $s_v=s_a$.

% 中文：预测和 Slack 更新在 Controller 中异步执行，调用调度和 token 生成无需等待。
Prediction and slack updates run asynchronously in the controller, so
invocation dispatch and token generation do not wait for them.

\subsection{Router and Connector}
\label{sec:router}

% 中文：Router 根据队列负载和前缀 KV 的位置调度调用，与此同时 Router 还会决定如果 Failure 发生时，该 Agent 调用 retry 到哪个DP Replica
% 这一选择也决定了如果要备份该调用对应的 KV，备份在哪里。
% retry 目标位于执行节点的受保护故障域之外，选择时考虑负载、可用内存和复制路径竞争。
Router schedules invocations based on queue load and prefix locality.
It also selects a retry DP replica for each Agent invocation in case of
failure. If the invocation's KV is selected for backup, it is placed at
this retry target. The target lies outside the execution worker's
protected failure domain and is chosen based on load, available memory,
and copy-path contention.

% 中文：在该目标上，主机内存（Peer DRAM）保存受保护 Agent 的完整 KV 检查点；GPU 内存
% （Peer HBM）额外保存 KV Master 从其前 L 层中选出的 KV 块。容量不足时，Router 将缺口
% 返回 KV Master。
At this target, host memory (\emph{Peer DRAM}) holds a protected Agent's
complete KV checkpoint, and GPU memory (\emph{Peer HBM}) additionally
holds the KV blocks that KV Master selects from its first $L$ layers.
If a placement does not fit, Router returns the capacity shortfall to
KV Master.

% 中文：各 TP rank 的 Connector 执行 KV 复制并向 KV Master 报告完成状态。
% 故障后 Router 使旧执行失效，并将调用重定向到 retry 目标；Engine 直接复用目标
% HBM 中的 KV，Connector 将 DRAM 中的 KV 搬入 GPU。目标成为执行节点后，
% Router 为后续备份重新选择 Peer。
Rank-local Connectors execute KV copies and report completion to KV
Master. After a failure, Router fences the old execution and redirects
the invocation to the retry target. The engine reuses resident Peer HBM
KV, while Connectors move DRAM-resident KV into GPU memory.
If that target becomes the execution worker, Router selects a new peer
for subsequent backups.

% 中文：上述路由与传输机制处理已选状态；接下来由 KV Master 决定保护哪些
% Agent 历史，以及保留其中哪些 KV 区域。
These routing and transfer mechanisms operate on selected state.
KV Master determines which Agent histories and which KV regions to
protect, as described next.

\subsection{KV Master}
\label{sec:master}

% 中文：KV Master 决定每个 retry 目标为哪些 Agent 保存什么。遵循 O1 和 O2，它在两个层次上
% 做决定：在 Agent 之间，只为 Slack 掩盖不了重试延迟的 Agent 在 Peer DRAM 中保存完整的
% KV 检查点；在受保护的 Agent 内部，再把恢复时需要等待的少量 KV 块放入 Peer HBM。两个层次
% 把覆盖与速度分开：DRAM 检查点保证恢复不必重建完整历史，HBM 副本只缩短恢复时间，可以随时
% 丢弃。下面先介绍如何选择 Agent，再介绍如何在受保护的 Agent 内部选择 KV，最后介绍副本管理与恢复。
KV Master decides what each retry target holds, and for which Agents.
Following O1 and O2, it decides at two levels. Across Agents, it keeps a
complete KV checkpoint in Peer DRAM only for Agents whose retry delay
their slack cannot absorb. Within a protected Agent, it also places in
Peer HBM the few KV blocks that recovery would otherwise wait for. The two
levels separate coverage from speed: the DRAM checkpoint ensures that
recovery never reconstructs the full history, while HBM copies only
shorten recovery and can be dropped at any time. We first describe how
Master selects Agents, then how it selects KV within a protected Agent,
and finally how it manages replicas and recovers.

\subsubsection{Agent Selection}
\label{sec:agent-protection}

% 中文：令 T_a 为通过完整重建（即对 Agent a 的历史做 Prefill）重试其被中断调用的延迟，
% Master 根据离线测得的 Prefill 延迟与历史长度的关系估计 T_a。将 R_i = T_a 和 Tracker 的
% Slack σ_i = s_a 代入式 (1)，得到 Slack 无法掩盖的重试延迟 E_a。
Let $T_a$ be the delay of retrying Agent $a$'s interrupted invocation by
full reconstruction, i.e., by prefilling its history. Master estimates
$T_a$ from an offline profile of prefill latency versus history length.
Applying Equation~\ref{eq:isolated_penalty} with $R_i=T_a$ and Tracker's
guarded slack $\sigma_i=s_a$ gives the retry delay that slack cannot hide,
\begin{equation}
 E_a=\bigl[T_a-s_a\bigr]^+.
 \label{eq:exposed-delay}
\end{equation}

% 中文：当 E_a > λ 时，Master 保护 a，即在 Peer DRAM 中保存其完整 KV 检查点。λ 是每次
% 故障可容忍的工作流延迟：未受保护的 Agent 重建时最多给工作流增加 λ；λ = 0 时，重建时间
% 超过 Slack 的 Agent 都受保护。若 Peer DRAM 放不下所有超过 λ 的 Agent，Master 按 E_a
% 从大到小接纳。T_a 随历史增长，s_a 随工作流推进变化，因此 Master 在调用边界和 Tracker
% 更新 Slack 时重新判断；降到 λ 以下的 Agent 释放其检查点。
Master protects $a$, i.e., keeps a complete KV checkpoint of $a$ in Peer
DRAM, when $E_a>\lambda$. The threshold $\lambda$ is the workflow delay
that the operator tolerates per failure: reconstructing an unprotected
Agent adds at most $\lambda$ to its workflow, and $\lambda=0$ protects
every Agent whose reconstruction exceeds its slack. If Peer DRAM cannot
hold every Agent above $\lambda$, Master admits them in decreasing $E_a$.
Both terms change during execution, as $T_a$ grows with the history and
$s_a$ changes as the workflow advances. Master therefore re-evaluates the
rule at invocation boundaries and whenever Tracker updates a slack, and an
Agent that falls below $\lambda$ releases its checkpoint.

% 中文：图 X(a) 将该规则用于图 2(b) 中的 Agent。A1 的重建落在 Slack 内；取 λ = 17 s 时，
% A2 和 A4 获得 DRAM 检查点，而 A3 的 16.3 s 不超过 λ，失败后重建。
Figure~\ref{fig:master-selection}(a) applies this rule to the Agents of
Figure~\ref{fig:agent_value}(b). A1's reconstruction fits within its
slack. With $\lambda=17$\,s, A2 and A4 receive DRAM checkpoints, whereas
A3's 16.3\,s does not exceed $\lambda$, so A3 is reconstructed if it fails.

\paragraph{DRAM checkpoints.}
\label{sec:master-checkpoint}
% 中文：受保护 Agent 的检查点按固定的 Z 个 token 推进，本文取 Z = 256，即 4 个 64-token
% 块。每新接受 Z 个 token，Connector 将这一 chunk 所有层的 KV 复制到 Peer DRAM；所有分片
% 到达后，Master 发布新的检查点位置。故障后，retry 目标恢复检查点，并重建其后至多 Z 个已接受
% token 的 KV，与 checkpoint-and-replay 恢复一致。Z 在正常执行开销与恢复工作之间折中：Z 越小，
% 传输越碎、发布越频繁；Z 越大，需要重建的尾部越长。Z = 256 时，重建只是一次很短的 Prefill，
% 而每个 token 的 KV 只以大块传输复制一次。
A protected Agent's checkpoint advances in fixed chunks of $Z$ tokens; we
use $Z=256$, i.e., four 64-token blocks. Once $Z$ new tokens are
accepted, Connectors copy the chunk's KV for all layers to Peer DRAM, and
Master publishes the new checkpoint position after all shards arrive.
After a failure, the retry target restores the checkpoint and rebuilds
the KV of the at most $Z$ accepted tokens after it, as in
checkpoint-and-replay recovery~\cite{strati2024dejavu,cao2026lumen}. $Z$
trades normal-execution cost against recovery work. A smaller $Z$ splits
copies into smaller transfers and publishes more often, whereas a larger
$Z$ lengthens the rebuilt tail. With $Z=256$, the rebuild is one short
prefill, and each token's KV is copied once, in a large transfer.

% 中文（图示）：(a) 使用图 2(b) 的未掩盖延迟；(b) 的访问计数与块布局为示意。
\begin{figure}[t]
 \centering
 \includegraphics[width=\columnwidth]{figures/master-selection.pdf}
 \caption{Selection by KV Master. (a) Retry delays not hidden by slack for
 the Agents of Figure~\ref{fig:agent_value}(b); with $\lambda=17$\,s, A2
 and A4 receive complete DRAM checkpoints. (b) Access counts of one
 protected Agent's blocks. Every checkpointed block is in DRAM; with $K=2$
 and $L=2$ (dashed line), the two most counted blocks of L1 and L2 are also
 in HBM, and the tail of at most $Z$ tokens is rebuilt during recovery.}
 \label{fig:master-selection}
 \Description{Top: horizontal bars of unhidden retry delay for Agents A1
 to A4, with 0, 32.6, 16.3 and 18.0 seconds, and a dashed threshold line
 at 17 seconds; A2 and A4 exceed it. Bottom: a grid of access counts for
 layers L1 to L4 and blocks B0 to B5 with a hatched tail column. In L1 the
 counts 9 and 7 of B0 and B3, and in L2 the counts 8 and 6 of B1 and B3,
 are highlighted as HBM blocks; a dashed line below L2 marks depth L.}
\end{figure}

\subsubsection{KV Selection}
\label{sec:region-selection}

% 中文：DRAM 检查点使受保护的 Agent 不必完整重建，但从 DRAM 恢复仍可能让调用等待传输。
% 从 DRAM 恢复可以与重建重叠：重建逐层进行，GPU 计算前面的层时，Connector 可以
% 传输后面层的 KV（图 3(b)）。但仍有两部分传输暴露：第 1 层之前没有计算可以掩盖其传输；
% 当一层 KV 的传输慢于一层的重建时，等待逐层累积。Master 把受保护 Agent 前 L 层中每层的
% Top-K 块同时放入 Peer HBM，以消除这两部分等待。这需要回答三个问题：选哪些块、保存多少层、
% 以及 HBM 放不下时优先保护哪些 Agent。
A DRAM checkpoint spares a protected Agent full reconstruction, but its
recovery may still wait for DRAM transfers. Restoring from DRAM overlaps
with the rebuild. The rebuild proceeds layer by layer, so Connectors load a later layer's KV while the GPU computes
earlier layers (Figure~\ref{fig:kv_value}(b)). Two parts of the transfer
nonetheless remain exposed. No computation precedes layer~1 to hide its
transfer, and when a layer's KV loads more slowly than a layer rebuilds,
waiting accumulates across layers. Master removes both by placing in Peer
HBM the Top-$K$ blocks of each of the first $L$ layers of a protected
Agent. This raises three questions: which blocks to place, how many layers
to cover, and which Agents to favor when HBM is short.

\paragraph{Top-$K$ blocks.}
% 中文：O2 表明稀疏注意力的读取集中在少数块上，Master 因此按访问次数选块。稀疏注意力在每层
% 为每个 query token 原生地选出要读取的 KV 块。引擎为第 ℓ 层的每个块 b 维护计数 c_{ℓ,b}：
% 一个 query 在任一 head 上选中 b 即计一次，与图 3(a) 的统计口径一致。每层计数最大的 K 个块
% 构成该层的 Top-K 集合。K 按模型固定，由离线测得的访问集中度确定。图 5(b) 中，K = 2、
% L = 2，L1 的 B0、B3 和 L2 的 B1、B3 同时放入 HBM，其余已检查点的块只在 DRAM 中。
O2 shows that sparse attention concentrates its reads on a few blocks, so
Master selects blocks by their access counts. At each layer, sparse
attention natively selects the KV blocks that each query token reads. The
engine keeps a count $c_{\ell,b}$ for each block $b$ of layer $\ell$ and
increments it once for each query that selects $b$ in any head, as in
Figure~\ref{fig:kv_value}(a). The $K$ blocks with the largest counts form
layer $\ell$'s Top-$K$ set. $K$ is fixed per model, from its access
concentration measured offline. In Figure~\ref{fig:master-selection}(b),
$K=2$ and $L=2$, so B0 and B3 of L1 and B1 and B3 of L2 are also in HBM,
while the other checkpointed blocks are only in DRAM.

\paragraph{Updating Top-$K$ sets.}
% 中文：Top-K 集合随检查点更新，而不是逐 token 更新。逐 token 的工作只有计数：注意力 kernel
% 为一个 query 选完块后，引擎对被选块的计数器做一次 scatter-add；这些索引本来就存在，操作在
% GPU 上完成，不需要与主机同步。每个检查点（即每 Z 个 token）Master 读取计数，乘以衰减因子
% β（0 < β < 1）使近期访问占主导，然后对前 L 层重新排序，新 chunk 的块也参与排序。为避免计数
% 相近的块来回替换，一个块只有在计数超过集合中最小成员的 (1 + ε) 倍时才进入 Top-K。进入的块
% 由 peer 从自己的 DRAM 复制到 HBM，不经过网络；离开的块在正在进行的读取结束后释放。所有层的
% 计数随检查点一起保存，后续层只在恢复需要时才排序。由于更新跟随检查点，HBM 中的块总在 DRAM
% 检查点中。
Master updates the Top-$K$ sets once per checkpoint rather than per
token. Per-token work is limited to counting: after the attention kernel
selects blocks for a query, the engine increments the counters of the
selected blocks with one scatter-add over indices that already exist, on
the GPU and without host synchronization. At each checkpoint, i.e., every
$Z$ tokens, Master reads the counts, multiplies them by a decay factor
$0<\beta<1$ so that recent accesses dominate, and re-ranks each of the
first $L$ layers together with the blocks of the new chunk. To keep blocks
with similar counts from alternating, a block enters a Top-$K$ set only if
its count exceeds that of the set's smallest member by a factor $1+\epsilon$.
The peer copies entering blocks from its own DRAM to its HBM, without
crossing the network, and releases leaving blocks once their active reads
finish. The counts of all layers are saved with the checkpoint, so later
layers are ranked only when a recovery needs them. Because updates follow
checkpoints, every HBM block is also in the DRAM checkpoint.

\paragraph{Layer-ordered recovery with misses.}
% 中文：恢复时，引擎逐层重建尾部至多 Z 个 token，稀疏注意力照常为尾部 token 选块。选块读取
% 的块级元数据（如压缩键）很小，属于检查点的一部分，前 L 层的元数据与 Top-K 块一起放在 HBM
% 中。对后续层，Connector 按层序先传元数据和 Top-K 块，再传其余 KV。选中的块若已在 GPU 中
% 即为命中；否则为缺失（miss），Connector 立即从 DRAM 召回，优先于按层序的传输。Kernel 在
% 缺失块传输期间先计算命中块上的注意力，再按 log-sum-exp 合并两部分结果，因此该层只等待
% 缺失块。由于 Top-K 覆盖了大部分访问（O2），且尾部至多 Z 个 token，缺失块很少；但它们只有
% 在所在层开始时才可知，无法预取。
During recovery, the engine rebuilds the tail of at most $Z$ tokens layer
by layer, and sparse attention selects blocks for the tail tokens as
usual. The block metadata that selection reads, such as compressed keys,
is small; it is part of the checkpoint, and HBM holds it for the first $L$
layers together with their Top-$K$ blocks. For later layers, Connectors
load the metadata and Top-$K$ blocks first in layer order and the rest of
the KV afterwards. A selected block already in GPU memory is a hit. Any
other selected block is a miss, which Connectors fetch from DRAM
immediately, ahead of the layer-ordered stream. The kernel computes
attention over the hits while the misses load and merges the two partial
results by their log-sum-exp, so the layer waits only for its misses.
Misses are few, because Top-$K$ sets cover most accesses (O2) and the tail
spans at most $Z$ tokens, but they become known only when their layer
starts and therefore cannot be prefetched.

% 中文：图 X 示意这一过程。(a) 先从 DRAM 传完全部 KV 再重建，整个传输都暴露，但此后没有
% 缺失。(b) 按层序先传 Top-K 块，传输与重建重叠；但每层开始时召回的缺失块插到流的前面，推迟
% 后续层的 Top-K 块，当 Top-K 块与缺失块的传输慢于重建时，每层都要等待。(c) L1–L4 的 Top-K
% 块在 HBM 中，重建立即开始；缺失块在命中块的计算期间到达，后续 Top-K 块也在需要之前到达，
% 不再等待；其余 KV 在调用就绪后继续传输。
Figure~\ref{fig:recovery-pipeline} illustrates this process. Loading the
whole KV from DRAM before the rebuild exposes the entire transfer but
leaves no misses (a). Loading Top-$K$ blocks first in layer order overlaps
the transfer with the rebuild (b). The misses of each layer, however, are
fetched when the layer starts and delay the Top-$K$ blocks of later
layers, so every layer waits when Top-$K$ blocks and misses load more
slowly than a layer rebuilds. With the Top-$K$ blocks of L1--L4 in HBM,
the rebuild starts at once, each layer's misses arrive while it computes
on hits, and later Top-$K$ blocks arrive before they are needed, so no
layer waits (c). The rest of the KV continues to load after the
invocation is ready.

% 中文（图示）：示意图，不对应特定模型或实测：8 层，时间以一层重建时间为单位；Top-K 块占
% 每层 KV 的一半，传输需 1.2 个单位，每层缺失块的召回需 0.3 个单位，即带宽低于 B*。
\begin{figure}[t]
 \centering
 \includegraphics[width=\columnwidth]{figures/recovery-pipeline.pdf}
 \caption{Illustrative recovery of one protected Agent with eight layers. Each
 dark segment rebuilds one layer for the tail of at most $Z$ tokens. Top-$K$
 blocks are half of each layer's KV; together with the misses fetched when
 each layer starts, they load more slowly than a layer rebuilds. (a) The
 whole KV loads from DRAM before the rebuild. (b) Top-$K$ blocks load first
 in layer order. (c) The Top-$K$ blocks of L1--L4 are in HBM. The rest of
 the KV loads after the Top-$K$ blocks.}
 \label{fig:recovery-pipeline}
 \Description{Three two-lane timelines, H2D above GPU. (a) The GPU waits
 while all KV loads, then rebuilds eight layers. (b) Top-K blocks load
 layer by layer, a short miss fetch is inserted whenever a layer starts,
 and every layer waits before it rebuilds; the rest of the KV loads
 afterwards. (c) With the first four layers' Top-K blocks in HBM, the eight
 layers rebuild without waiting while Top-K blocks and miss fetches
 interleave on H2D.}
\end{figure}

\paragraph{Depth $L$.}
% 中文：设模型有 N 层，r 为一层尾部重建的时间，h 为在有效 H2D 带宽 B 下传输一层 Top-K 块与
% 缺失块的时间；B* 为 h = r 时的带宽。前 L 层在 HBM 中时，第 ℓ > L 层最早在 (ℓ−1)r 开始，
% 其 Top-K 块在 (ℓ−L)h 到达。任何层都不等待当且仅当对所有 ℓ ≤ N 都有 (ℓ−L)h ≤ (ℓ−1)r，
% 即式 (3)。这里把缺失块近似计入每层的传输量。第 1 层之前没有计算可以掩盖传输，因此 L ≥ 1。
Let the model have $N$ layers, let $r$ be the time to rebuild one layer
for the tail, and let $h$ be the time to load one layer's Top-$K$ blocks
and misses at effective host-to-device bandwidth $B$; $B^*$ is the
bandwidth at which $h=r$. With the first $L$ layers in HBM, layer
$\ell>L$ can start at $(\ell-1)r$, and its Top-$K$ blocks arrive by
$(\ell-L)h$. No layer waits if $(\ell-L)h\le(\ell-1)r$ for every
$\ell\le N$. Since nothing hides layer~1's transfer, this gives
\begin{equation}
 L=\max\Bigl\{1,\;\Bigl\lceil N-(N-1)\,\frac{B}{B^*}\Bigr\rceil\Bigr\},
 \label{eq:depth}
\end{equation}
where misses are counted as part of each layer's load.

% 中文：图 Y 绘出式 (3)。带宽不低于 B* 时 L = 1；低于 B* 时 L 随带宽下降线性增大。B 是单个
% 恢复实际获得的带宽：传输很少能打满链路，并发恢复还会共享链路，因此 B 低于链路峰值，需要在
% 预期的并发恢复下离线测量。部署据此为每个模型设定 L。
Figure~\ref{fig:depth-bandwidth} plots Equation~\ref{eq:depth}. At or
above $B^*$, $L=1$. Below $B^*$, $L$ grows linearly as the bandwidth
drops, and without HBM the waiting grows quickly. $B$ is the bandwidth
that one recovery actually obtains. Transfers rarely saturate the link,
and concurrent recoveries share it, so $B$ lies below the link's peak and
is measured offline under the expected concurrent recoveries. The
deployment sets $L$ for each model from this profile.

% 中文（图示）：示意图，32 层，与图 X 采用相同的规则。
\begin{figure}[t]
 \centering
 \includegraphics[width=\columnwidth]{figures/depth-bandwidth.pdf}
 \caption{Illustrative depth $L$ versus effective host-to-device bandwidth
 for a 32-layer model. Bandwidth is normalized to $B^*$, at which a layer's
 Top-$K$ blocks and misses load in one layer's rebuild time. (a) Smallest
 $L$ for which no layer waits (Equation~\ref{eq:depth}). (b) Waiting
 without HBM, in layer rebuild times. Below $B^*$ (dashed), loads fall
 behind the rebuild and $L$ grows as the bandwidth drops.}
 \label{fig:depth-bandwidth}
 \Description{Two stacked panels over bandwidth from 0 to 2 times B-star.
 (a) Depth L falls in steps from 31 near zero bandwidth to 1 at B-star and
 stays 1 above. (b) On a logarithmic scale, waiting without HBM falls from
 several hundred layer rebuild times near zero bandwidth to one at B-star,
 then declines slowly.}
\end{figure}

\paragraph{HBM admission.}
% 中文：同一模型的每个受保护 Agent 占用相同的 HBM：L·K 个块加上这些层的元数据。令 T_a^D 和
% T_a^H 分别为仅从 DRAM 恢复和有 HBM 副本时的重试延迟，HBM 副本为工作流减少的延迟为 G_a。
% Master 按 G_a 从大到小把受保护 Agent 接纳到 Adapter 给出的配额中，并跳过 G_a = 0 的
% Agent，即 DRAM 恢复已被 Slack 掩盖的 Agent。由于各 Agent 占用相同，这一顺序在配额内减少的
% 延迟最多。放不下的 Agent 保留 DRAM 检查点，Master 将其报告给 Adapter，作为增大配额的信号。
Every protected Agent of a model occupies the same HBM: $L\cdot K$ blocks
plus the metadata of these layers. Let $T^{D}_a$ and $T^{H}_a$ be $a$'s
retry delays when recovering from DRAM alone and with its HBM copy. The
HBM copy removes
\begin{equation}
 G_a=\bigl[T^{D}_a-s_a\bigr]^+-\bigl[T^{H}_a-s_a\bigr]^+
 \label{eq:hbm-gain}
\end{equation}
of workflow delay. Master admits protected Agents to Adapter's HBM quota
in decreasing $G_a$ and skips Agents with $G_a=0$, whose recovery from
DRAM already fits within their slack. Because all copies have the same
size, this order removes the most delay within the quota. Agents that do
not fit keep their DRAM checkpoint, and Master reports them to Adapter as
its signal for growing the quota (Section~\ref{sec:capacity}).

\subsubsection{Replica Management and Recovery}
\label{sec:master-online}
\label{sec:execution}

% 中文：副本的生命周期跟随 Agent 而不是调用。同一 Agent 的后续调用扩展同一前缀，因此副本跨
% 调用有效；Agent 完成任务或降到 λ 以下时，Master 释放其副本。复制是异步的：副本目录记录每个
% 块的层、token 区间、TP 分片和位置，所有分片到达后才发布检查点或 HBM 副本。
Replicas follow Agents rather than invocations. Later invocations of an
Agent extend the same prefix, so its replicas remain valid across
invocations, and Master releases them when the Agent finishes its task or
falls below $\lambda$. Because copies are asynchronous, a replica manifest
records each block's layer, token range, TP shards, and location, and
publishes a checkpoint or HBM copy only after all its shards arrive.

% 中文：故障后，Router 屏蔽旧执行并将被中断的调用重定向到 retry 目标。目标从副本目录中找到
% 最近发布的检查点，引擎直接使用 HBM 中的块和元数据，Connector 按上述方式从 DRAM 加载其余
% KV 和缺失块，引擎逐层重建检查点之后的尾部，然后继续生成。DRAM 检查点是完整的，因此 HBM
% 副本只影响恢复时间，不影响正确性；未发布的复制在恢复中被忽略。
After a failure, Router fences the old execution and redirects the
interrupted invocation to its retry target. The target looks up the
latest published checkpoint in the manifest. The engine uses the HBM
blocks and metadata in place, Connectors load the remaining KV and misses
from DRAM as described above, and the engine rebuilds the tail layer by
layer before resuming generation. Because the DRAM checkpoint is complete,
HBM copies affect recovery time but never correctness, and recovery
ignores unpublished copies.

% 中文：正常执行的开销与 token 数成比例的部分只有计数。未受保护的 Agent 没有开销；受保护的
% Agent 每个 token 只做计数器自增。每个检查点，Connector 复制一个 Z-token chunk 到 Peer
% DRAM，Master 对 L 层重新排序，peer 在本地把至多 L·K 个进入 Top-K 的块从 DRAM 复制到
% HBM；上述滞后条件使这一复制只针对计数明显变化的块。
Normal execution pays little, and mostly once per $Z$ tokens.
Unprotected Agents incur no cost. A protected Agent adds only counter
increments per token. Per checkpoint, Connectors copy one $Z$-token chunk
to Peer DRAM, Master re-ranks $L$ layers, and the peer locally copies at
most $L\cdot K$ entering blocks from its DRAM to its HBM; the hysteresis
restricts these copies to blocks whose counts change substantially.

\subsection{Adapter}
\label{sec:capacity}

% 中文：O3 表明前台服务可让出的 HBM 随负载变化。Adapter 为每个接管 rank k 维护 HBM 副本
% 配额 B_k，KV Master 只在该配额内接纳 Agent（每个 Agent 占 L·K 个块）。配额受上限 C_k
% 约束：B^pool_k 为扣除权重和固定分配后的显存，B^serve_k 为服务高水位与执行预留，
% B^reserve_k 为增长与 workspace 余量，B^cap_k 为管理员设定的上限。
O3 shows that the HBM that foreground serving can spare changes with
load. Adapter maintains an HBM replica quota $B_k$ for each takeover rank
$k$, and KV Master admits Agents, each occupying $L\cdot K$ blocks, only
within it. The quota is bounded by
\begin{equation}
 C_k=\max\{0,\min(B^{\rm cap}_k,
 B^{\rm pool}_k-B^{\rm serve}_k-B^{\rm reserve}_k)\},
 \label{eq:ceiling}
\end{equation}
where $B^{\rm pool}_k$ is the memory left after weights and fixed
allocations, $B^{\rm serve}_k$ covers the serving high watermark and
execution reservations, $B^{\rm reserve}_k$ keeps headroom for growth and
workspace, and $B^{\rm cap}_k$ is an operator limit.

% 中文：Adapter 以 AIMD 调整 B_k。分配卡顿、预留被侵占或因显存不足的准入等待出现时，
% B_k 乘以 γ < 1；连续 m 个窗口稳定且 Master 报告有 Agent 因 HBM 不足未被接纳时，B_k
% 增加 Δ，并截断到 C_k。深度 L 与 K 不随配额变化，配额只改变可放入 HBM 的 Agent 数量。
Adapter adjusts $B_k$ by additive increase and multiplicative
decrease~\cite{chiu1989aimd}. Allocation stalls, reserve violations, or
admission waits caused by memory multiply $B_k$ by $\gamma<1$. After $m$
stable windows, Adapter adds $\Delta$ if Master reports protected Agents
deferred for HBM capacity, and clips the result to $C_k$. $L$ and $K$ stay
fixed, so the quota changes only how many Agents HBM holds.

% 中文：配额缩小时，Master 停止新的 HBM 接纳，并按 G_a 从小到大释放 Agent
% 的 HBM 副本，等其读取结束后回收。DRAM 检查点是完整的，因此无需迁移数据，这些 Agent
% 只是改为从 DRAM 恢复。前台准入按实际占用而非配额判断，释放完成前不会超额分配。
When the quota shrinks, Master stops new HBM admissions and releases the
HBM copies of Agents in increasing $G_a$ once their readers
drain. Because their DRAM checkpoints are complete, no data moves, and
these Agents recover from DRAM instead. Foreground admission follows
actual occupancy rather than the quota, so memory is not overcommitted
before the release completes.
